\section{Purpose and status of the proposed approach}
\label{sec:bp_scope}

The compact formulation assigns flights and maintenance decisions directly to
aircraft. Chapter~\ref{ch:strengthening} reduces its size and strengthens some
of its inequalities. A complementary approach is to optimize over complete
integer scheduling patterns for individual tails, generating those patterns
only when they can improve a coordinating master problem. This chapter develops
that approach for the full A/B/C/D maintenance hierarchy through
Dantzig--Wolfe decomposition and branch-and-price.

\textit{This chapter is a methodological proposal submitted for review before
implementation. No branch-and-price solver, measured speedup, generated root
bound, or computational optimality result is claimed here. The results in
Chapter~\ref{ch:experiments} concern the previously implemented methods.}

The target is an exact reformulation of the selected compact model, not a new
interpretation of maintenance legality. The initial proposal retains the
aircraft-indexed binary trigger formulation, hard flight coverage, full ABCD
checks, and the same assignment domains, initial conditions, hierarchy,
maintenance durations, deferral policy, cost accounting and routing flags.
The optional global-trigger formulation is outside this initial decomposition
because its cross-aircraft variables require a separate separability analysis.
In particular, replacing a compact constraint by a stricter physical replay
rule would change the comparison problem and is not part of this proposal.

For an unambiguous reference, the equivalence claim in this chapter concerns
the default full-maintenance profile of \texttt{MILP\_Sheduler}: hard coverage;
aircraft-indexed sparse $z$ variables; day-specific $y$ and $\gamma$ bounds;
C2--C3 and interval-clique overlap enabled; capacity, hierarchy, C12, C12b,
C13b and corrected split C13 enabled; three-day A/B deferral; global
$M_{\mathrm{BIG}}=9{,}999{,}999$ for C13/C13b; and the default weak C9 link.
The C14b duration constant is $M_{\mathrm{C14b}}=\max_c K_c^{\mathrm{dur}}+1$.
The global-trigger, soft-coverage, no-capacity, reachability, tight-C13-$M$,
original-paper-C13, strong-link, strong-conflict and sanity-fixing options are
disabled. The documented compatibility C15 component is empty; C8 supplies
the maintenance-flight blocking rows. Any experiment with another option set
requires the corresponding local set and coefficient definitions to change.

\section{Literature review and positioning}
\label{sec:bp_literature}

\subsection{Decomposition and exact integer optimization}

\citet{dantzig1960decomposition} distinguish constraints belonging to individual
subproblems from constraints connecting those subproblems. Their coordinating
linear program uses columns proposed by the subproblems, while its dual prices
define successive subproblem objectives. The relevance here is the separation
between a tail's routing and maintenance decisions and the fleet-wide demands
for flight coverage and maintenance-station capacity.

Decomposing an integer problem requires more than solving the coordinating LP.
\citet{barnhart1998branchprice} describe column generation within a
branch-and-bound tree, including pricing at tree nodes and branching rules
compatible with implicit columns. This distinction is essential: column
generation at the root yields a relaxation bound, whereas exact integer
optimization requires pricing under the restrictions imposed throughout the
search tree. An integer solve over a fixed, generated column pool can produce
an incumbent, but does not certify the optimum over omitted columns.

\subsection{Airline strings, routing and maintenance}

\citet{barnhart1998flight} formulate aircraft fleeting and routing with flight
strings, allowing maintenance and connection requirements to be represented
within a route-based description. That work motivates replacing individual
assignment decisions by structured aircraft duties. It is distinct from the
general branch-and-price paper by \citet{barnhart1998branchprice}, although both
were published in 1998.

\citet{desaulniers1997daily} develop daily routing models using
set partitioning and time-constrained networks, with column
generation and compatible branching. Their profit-maximization setting must
not be confused with the minimization problem in this thesis: the direction
of a relaxation bound depends on the objective sense. Their methodological
relevance is the connection between decomposed routing formulations and
branching decisions that remain expressible in pricing.

\citet{sarac2006branch} specifically address operational aircraft maintenance
routing through branch-and-price. This provides direct application precedent
for generating maintenance-aware aircraft patterns. It does not by itself
establish that their maintenance conventions equal the ABCD hierarchy,
initial-state treatment and station-capacity accounting used here. Those
model-specific details must be stated and validated independently.

The compact model of \citet{khaled2018compact} remains the formulation reference
for the thesis, together with the modifications in
Chapters~\ref{ch:maintenance} and~\ref{ch:hierarchy}. The proposed contribution
is therefore not the invention of branch-and-price. It is a precisely defined
full-ABCD decomposition of that reference model, with reconstructable pattern
certificates, transparent lower-bound certification and a controlled evaluation
against the compact formulation.

\begin{table}[htbp]
\centering\small
\caption{Distinctions between route-based solution procedures.}
\label{tab:bp_method_distinctions}
\begin{tabularx}{\textwidth}{@{}p{0.25\textwidth}XX@{}}
\toprule
Procedure & Columns considered & What can be certified \\
\midrule
Complete pattern enumeration & Every feasible local integer pattern & Full master LP and integer optima, if solved \\
Root column generation & Improving patterns generated at the root & Full master LP bound after exact pricing certification \\
Restricted integer master & A fixed generated pool & A feasible incumbent; not generally full integer optimality \\
Branch-and-price & Patterns generated under restrictions at every node & Integer optimality if the complete search is certified \\
\bottomrule
\end{tabularx}
\end{table}

\section{Tail-local integer scheduling patterns}
\label{sec:bp_local_sets}

Flights, tails, stations, days and check types retain the sets
$\mathcal F$, $\mathcal P$, $\mathcal A$, $\mathcal D$ and
$\mathcal C=\{A,B,C,D\}$, respectively. Let $\mathcal K\subseteq\mathcal A$
be the maintenance stations. For each tail $j$, collect all of its original
binary decisions into a vector $v_j$. Its components include $x_{ij}$,
$z_{ijdc}$, $y_{jdc}$, $\gamma_{jdc}$ and maintenance-start auxiliaries.
The mathematical symbol $\gamma_{jdc}$ corresponds to \texttt{mega[j,d,c]}
in the implementation. Variables absent from a sparse domain are represented
as zero when coefficients are evaluated.

The selected compact model's local feasible set is
\begin{equation}
 \mathcal Q_j=
 \left\{v_j\in\{0,1\}^{n_j}:
 A_jv_j\le h_j,\ E_jv_j=e_j,\ \ell_j^{\mathrm{var}}\le v_j\le u_j^{\mathrm{var}}
 \right\}.
 \label{eq:bp_local_set}
\end{equation}
The matrices denote exactly the previously defined tail-local constraints;
they do not introduce a replacement maintenance model. Table~\ref{tab:bp_split}
specifies their ownership. Flight coverage C1 and shared capacity C10 are
excluded from $\mathcal Q_j$ and retained in the master.

\begin{table}[htbp]
\centering\small
\caption{Constraint ownership in the proposed decomposition.}
\label{tab:bp_split}
\begin{tabularx}{\textwidth}{@{}p{0.19\textwidth}p{0.17\textwidth}X@{}}
\toprule
Constraint group & Owner & Meaning \\
\midrule
C1 & Master & Hard exact coverage of every flight \\
C10 & Master & Shared station/day capacity across tails and check types \\
C2--C3 & Tail pricing & Per-aircraft equipment flow and minimum turn \\
Overlap & Tail pricing & Per-aircraft temporal clique rows \\
C8 & Tail pricing & Same-day blocking and deferred-trigger escape blocking \\
C9 & Tail pricing & Default weak trigger/assignment link, exactly as coded \\
C11 & Tail pricing & Per-tail/day/check sum of triggers equals $y_{jdc}$ \\
Hierarchy & Tail pricing & $\gamma$ equalities from the selected checks \\
C12, C12b & Tail pricing & Calendar spacing and initial calendar deadline \\
C13, C13b & Tail pricing & Split hour rows and pre-horizon accumulated hours \\
C14, C14b & Tail pricing & At most one check/day and multi-day duration \\
C15 & No rows & Empty compatibility component; C8 performs blocking \\
Start links & Tail pricing & $s_{jdc}$ detects $y_{jdc}$ starts for objective cost \\
Variable domains & Tail pricing & Sparse $X,Z$, binary domains and day-specific zero bounds \\
Objective & Tail pricing & Assignment cost plus maintenance-start penalty \\
\bottomrule
\end{tabularx}
\end{table}

The table is the decomposition claim, row by row: every nonempty original
row other than C1 and C10 is copied without changing its index sets,
coefficients, sense or bounds into exactly one tail block. For example, C9
uses the implementation's default weak link: for C/D checks it links the
assignment to the first candidate trigger day, not the stronger sum-over-days
link. The sparse $Z$ domain, impossible-state upper bounds and hierarchy bounds
are also part of that block, not merely optional presolve. C15 has no active
rows, so it contributes nothing to either block.

For clarity, let $H_{jdd'}^c$ be the assigned flight-minute accumulation
in a C13 interval, let $Y_{jdd'}^c$ be the sum of intervening hierarchy reset
indicators, and let $T_c=60H_c$ be the flight-hour threshold in minutes when
the input threshold $H_c$ is expressed in hours. Under the selected default
profile, C13 uses the single global constant $M_{\mathrm{BIG}}=9{,}999{,}999$,
not the optional interval-specific tightening. The reference split rows are
\begin{align}
 H_{jdd'}^c&=\sum_{i\in\mathcal F_{d,d'}:(i,j)\in X}t_i x_{ij},\\
 Y_{jdd'}^c&=\sum_{r\in\mathcal D:\ d<r<d'}\gamma_{jrc}.
 \label{eq:bp_c13_accumulation}
\end{align}
where $\mathcal F_{d,d'}$ uses the same departure-day interval convention as
the compact builder. The enforced pair is
\begin{align}
 H_{jdd'}^c &\le T_c+M_{\mathrm{BIG}}Y_{jdd'}^c
                         +M_{\mathrm{BIG}}\gamma_{jdc},\\
 H_{jdd'}^c &\le T_c+M_{\mathrm{BIG}}Y_{jdd'}^c
                         +M_{\mathrm{BIG}}\gamma_{jd'c}.
 \label{eq:bp_retained_c13}
\end{align}
All terms use minutes. For each endpoint day $d'$, C13b is the additional
opening-horizon row
\begin{align}
 H_{jd'}^{c,0}&=\sum_{i\in\mathcal F_{0,d'}:(i,j)\in X}t_i x_{ij},\\
 H_{jd'}^{c,0}&\le 60(H_c-h_{jc}^{0})+
 M_{\mathrm{BIG}}\left(\sum_{r\in\mathcal D:r<d'}\gamma_{jrc}
                        +\gamma_{jd'c}\right),
 \label{eq:bp_c13b}
\end{align}
where $h_{jc}^{0}$ is the input's initial accumulated A/B flight-hour
exposure. The interval domains and opening-horizon rows are inherited from
the reference model. Branch-and-price does not, by
itself, prove these modelling conventions physically correct. Its equivalence
claims below concern the chosen mathematical model. Enabling
\texttt{use\_tight\_c13\_m} changes the LP relaxation and therefore defines a
different bound-comparison protocol, even when its interval constants are
valid for the same integer schedules.

The three-day A/B deferral cap and unrestricted \texttt{none} option define
different $\mathcal Q_j$ sets and must be compared separately. Likewise, the
compact flag that activates C2--C3 must not be interpreted as adding the
heuristic scheduler's explicit ferry events. The proposed initial protocol
keeps restrictive station-wait filtering disabled.

Let $\mathcal R_j$ index all elements of the finite set $\mathcal Q_j$, with
$v_j^r$ denoting pattern $r$. A pattern is a complete vector of decisions,
not merely a tuple of flight identifiers. Identical flight assignments with
different check days or trigger placements must remain available. An idle
pattern is admissible only if it satisfies every local constraint, including
initial C/D deadlines. A maintenance-only pattern may be used only if the
existing trigger formulation actually permits it.

\section{Pattern coefficients and the full master problem}
\label{sec:bp_master}

Define the coverage coefficient and full pattern cost by
\begin{align}
 a_{ijr}&=x_{ij}^r,\\
 C_{jr}&=\sum_{i\in\mathcal F:(i,j)\in X}c_{ij}x_{ij}^r
       +100\sum_{d\in\mathcal D}\sum_{c\in\mathcal C}
                    K_c^{\mathrm{dur}}s_{jdc}^r.
 \label{eq:bp_cost_coverage}
\end{align}
Here $s_{jdc}$ is the binary maintenance-start variable and
$K_c^{\mathrm{dur}}$ is the check duration in days. The compact objective
charges 100 times the full check duration for each start. The local start-link
rows include $s_{jdc}\le y_{jdc}$ and
$s_{jdc}\ge y_{jdc}-y_{j,d-1,c}$ (with $s_{jdc}\ge y_{jdc}$ on the first
horizon day). Thus a 0-to-1 transition forces $s=1$ and an inactive day forces
$s=0$; on a continuing 1-to-1 day either binary value satisfies these rows.
Because the objective coefficient is positive, an optimum chooses $s=0$ on
such continuation days. The exact local pattern set retains the original
variable and inequalities rather than tightening them. This term is part of
the current objective;
omitting it would make the master objective inequivalent even if every
feasible pattern were present. The assignment-cost lookup uses the actual
input flight row, not an assumption that identifiers are one-based. For any
candidate local vector $v_j$, write its complete objective contribution as
$f_j(v_j)$; thus $C_{jr}=f_j(v_j^r)$.

Let $\mathcal I_{dc}=\{i:d\in D(i,c)\}$ be the trigger flights returned by
the implementation's day/check index, and let
$\mathcal F_m(k,d)=\{i:\operatorname{dest}(i)=k,
\operatorname{arrivalDay}(i)\le d\}$. The exact trigger set for C10 is
$\mathcal T_{kdc}=\mathcal I_{dc}\cap\mathcal F_m(k,d)$. Only instantiated
aircraft/flight trigger indices are included. For a fixed-tail pattern $r$,
the capacity coefficient is
\begin{equation}
 b_{kdjr}=\sum_{c\in\mathcal C}\sum_{i\in\mathcal T_{kdc}}
                    \mathbf 1\{(i,j,d,c)\in Z\}z_{ijdc}^r.
 \label{eq:bp_capacity_coefficient}
\end{equation}
This reproduces C10's trigger/start accounting. It does not silently charge
capacity on every occupied day of a multi-day check. Local duration blocking
continues to be enforced by C14b and the other retained rows. Denote the
reference capacity by $K_{kd}$; when the input has a station-wide constant,
$K_{kd}=K_k$ for every day.

With $\lambda_{jr}=1$ selecting pattern $r$ for tail $j$, the full integer
master is
\begin{align}
 \min\quad &\sum_{j\in\mathcal P}\sum_{r\in\mathcal R_j}C_{jr}\lambda_{jr}
                  \label{eq:bp_master_obj}\\
 \text{s.t.}\quad
 &\sum_{j\in\mathcal P}\sum_{r\in\mathcal R_j}a_{ijr}\lambda_{jr}=1
                  &&\forall i\in\mathcal F,\label{eq:bp_master_cover}\\
 &\sum_{r\in\mathcal R_j}\lambda_{jr}=1
                  &&\forall j\in\mathcal P,\label{eq:bp_master_convexity}\\
 &\sum_{j\in\mathcal P}\sum_{r\in\mathcal R_j}b_{kdjr}\lambda_{jr}\le K_{kd}
                  &&\forall k\in\mathcal K,d\in\mathcal D,
                     \label{eq:bp_master_capacity}\\
 &\lambda_{jr}\in\{0,1\}
                  &&\forall j\in\mathcal P,r\in\mathcal R_j.
                     \label{eq:bp_master_binary}
\end{align}
The equality in \eqref{eq:bp_master_convexity} ensures a complete local state
is selected for every tail, even if its flight assignment is empty.

\begin{proposition}[Integer equivalence]
\label{prop:bp_integer_equivalence}
Suppose every local integer pattern of the selected compact model is included
in $\mathcal R_j$, and C1 and C10 are its only cross-tail constraints.
Then the full integer master and that compact model have equal optimal
objectives, with cost-preserving feasible-solution mappings in both directions.
\end{proposition}
\begin{proof}
A compact feasible solution determines one $v_j\in\mathcal Q_j$ for each tail.
Select its indexed column. C1 gives \eqref{eq:bp_master_cover}, C10 gives
\eqref{eq:bp_master_capacity}, and exactly one pattern is selected per tail.
Conversely, binary master variables and the convexity equations select one
member of $\mathcal Q_j$ per tail. Its local constraints hold, while the two
shared master groups reproduce C1 and C10. Equation~\eqref{eq:bp_cost_coverage}
preserves the objective in either mapping.
\end{proof}

The separability premise must be checked against the active implementation
options before this proposition is used computationally. The patterns are
not claimed to be physically connected routes beyond what the retained
compact routing constraints themselves establish.

\section{Continuous relaxation and bound strength}
\label{sec:bp_relaxation}

Replace \eqref{eq:bp_master_binary} by $\lambda_{jr}\ge0$. The convexity
equations already imply $\lambda_{jr}\le1$. Its optimum is denoted
$z_{\mathrm{DW}}$. A fractional selection projects into original variables as
\begin{equation}
 \bar v_j=\sum_{r\in\mathcal R_j}v_j^r\lambda_{jr}
              \in\operatorname{conv}(\mathcal Q_j).
 \label{eq:bp_projection}
\end{equation}
Let $P_j$ be the bound-preserving continuous relaxation of the same local
compact block. Because $P_j$ is convex and contains $\mathcal Q_j$,
$\operatorname{conv}(\mathcal Q_j)\subseteq P_j$.

\begin{proposition}[Relaxation ordering]
\label{prop:bp_bound_ordering}
For identical objective, local domains and shared constraints, and finite
optima of feasible models,
\begin{equation}
 z_{\mathrm{compact\,LP}}\le z_{\mathrm{DW}}\le z^*\le U,
 \label{eq:bp_bound_order}
\end{equation}
where $z^*$ is the common integer optimum and $U$ is any validated incumbent.
\end{proposition}
\begin{proof}
Every projected master-LP solution lies in each $P_j$ and satisfies the same
shared rows. It is therefore feasible in the compact relaxation, so its
minimum cannot be below the compact-LP minimum. Every integer master solution
is feasible in the master relaxation, giving the second inequality. The last
inequality is the definition of an incumbent upper bound.
\end{proof}

Equality is possible; no strict improvement or runtime advantage follows
from this proposition. The comparison uses the standard bound-preserving LP,
not a legacy procedure that reopens zero binary bounds. Even a stronger root
bound may be expensive to obtain if pricing requires difficult MILP solves.

\section{Restricted master, dual prices and exact pricing}
\label{sec:bp_pricing}

Since the full pattern sets can be enormous, begin with subsets
$\widehat{\mathcal R}_j\subseteq\mathcal R_j$. The restricted master problem
(RMP) is the master LP using only these subsets. In this section assume that
the RMP is feasible without artificial variables and solved to optimality.
Let $\pi_i$ and $\alpha_j$ be unrestricted dual variables for coverage and
convexity, respectively, and let $\beta_{kd}\le0$ be the dual variables for
capacity. The restricted dual is
\begin{align}
 \max\quad &\sum_{i\in\mathcal F}\pi_i+\sum_{j\in\mathcal P}\alpha_j
                  +\sum_{k\in\mathcal K,d\in\mathcal D}K_{kd}\beta_{kd}
                  \label{eq:bp_dual_obj}\\
 \text{s.t.}\quad
 &\sum_{i\in\mathcal F}a_{ijr}\pi_i+\alpha_j
                  +\sum_{k,d}b_{kdjr}\beta_{kd}\le C_{jr}
                  &&\forall j,r\in\widehat{\mathcal R}_j,\\
 &\pi_i,\alpha_j\in\mathbb R,\quad\beta_{kd}\le0.
\end{align}
The capacity sign is appropriate for a minimization primal with a
less-than-or-equal row. Thus $-\beta_{kd}b_{kdjr}$ penalizes use of a scarce
station-day resource in the reduced cost.

For each tail, solve the pricing MILP
\begin{align}
 p_j=\min_{v_j\in\mathcal Q_j}\quad
 &f_j(v_j)-\sum_{i\in\mathcal F}\pi_i x_{ij}
          -\sum_{k,d}\beta_{kd}b_{kd}(v_j)-\alpha_j.
 \label{eq:bp_pricing_milp}
\end{align}
Here $b_{kd}(v_j)$ is the linear trigger expression in
\eqref{eq:bp_capacity_coefficient}. A feasible pricing solution with negative
reduced cost defines a column to add; it need not be the pricing optimum to
be useful. However, failure of a heuristic to find such a column is not an
optimality certificate. Exact termination requires a proven lower bound on
the pricing objective for every tail.

The initial pricing design retains the compact local MILP rather than
replacing it by a resource-constrained shortest-path algorithm. A future
label-setting implementation would need an equivalence proof covering
alternative maintenance timings, hierarchy resets and initial conditions.
The existing greedy timeline simulator can suggest seed patterns only after
they have been checked against $\mathcal Q_j$. Its deterministic maintenance
placement does not enumerate every feasible column.

\section{Certified lower bounds with incomplete pricing}
\label{sec:bp_certification}

Restricting columns in a minimization LP can increase its objective. Hence
$z_{\mathrm{RMP}}\ge z_{\mathrm{DW}}$, but $z_{\mathrm{RMP}}$ alone need not
be a lower bound on $z^*$. A bound must account for columns that remain absent.

Let $\underline p_j\le p_j$ be a certified lower bound from pricing. A pricing
incumbent provides an upper bound on $p_j$, not the required lower bound.
Define $\delta_j=\min\{0,\underline p_j\}$.

\begin{proposition}[Pricing-corrected lower bound]
\label{prop:bp_corrected_bound}
With feasible optimal Phase-II RMP duals and a finite certified
$\underline p_j$ for every nonempty local pattern set,
\begin{equation}
 L=z_{\mathrm{RMP}}+\sum_{j\in\mathcal P}\delta_j
       \le z_{\mathrm{DW}}\le z^*.
 \label{eq:bp_certified_lb}
\end{equation}
If every pricing optimum is certified nonnegative, then
$z_{\mathrm{RMP}}=z_{\mathrm{DW}}$.
\end{proposition}
\begin{proof}
For every pattern of tail $j$, its reduced cost is at least
$p_j\ge\underline p_j\ge\delta_j$. Replace $\alpha_j$ by
$\alpha'_j=\alpha_j+\delta_j$ while keeping the other dual variables unchanged.
All reduced costs become the previous reduced costs minus $\delta_j$, hence
are nonnegative. These adjusted duals are feasible for every column of the
full master. Their objective is \eqref{eq:bp_certified_lb}, proving the bound
by weak duality. If the original duals already have no violated column,
they are feasible for the full dual; the RMP primal is also full-master
feasible, so their equal objectives prove optimality.
\end{proof}

The proof is stated in exact arithmetic. In implementation, solver feasibility
and objective-bound tolerances must be recorded and handled conservatively.
For example, if only $p_j\ge-\varepsilon$ is certified for every tail,
$z_{\mathrm{RMP}}-|\mathcal P|\varepsilon$ is the corresponding mathematical
lower bound; a zero correction is not justified by that tolerance alone.
Unresolved pricing with no valid finite bound yields no new certificate.
An inherited node bound can still be retained. A proven empty local set
implies node infeasibility, not a numeric pricing correction.

\section{Feasible initialization and Phase I}
\label{sec:bp_phase_one}

A small initial pool may not cover all flights or contain a pattern for every
tail. A feasible full integer solution is helpful but is not required to
initialize column generation. Introduce nonnegative artificial variables
$s_i^+,s_i^-$ for coverage and $t_j^+,t_j^-$ for convexity:
\begin{align}
 \sum_{j,r}a_{ijr}\lambda_{jr}+s_i^+-s_i^-&=1 &&\forall i,\\
 \sum_r\lambda_{jr}+t_j^+-t_j^-&=1 &&\forall j.
\end{align}
Retain station capacities and minimize the artificial violation
\begin{equation}
 \Phi=\sum_i w_i(s_i^++s_i^-)+\sum_j w_j(t_j^++t_j^-),
 \qquad w_i,w_j>0.
 \label{eq:bp_phase_one_obj}
\end{equation}
With nonnegative capacities, the all-zero real-column solution and positive
right-hand-side artificials provide an initial feasible Phase-I RMP. Real
columns have zero direct cost in Phase I. Their pricing objective uses
Phase-I duals and zero pattern cost, not $C_{jr}$ from Phase II.

If the artificial objective reaches zero, fix the artificial variables to
zero and optimize the original objective. If exact Phase-I pricing certifies
that the complete Phase-I optimum is positive, the full master LP is
infeasible, and therefore so is its integer restriction. Positive artificial
values in an incompletely priced RMP do not prove infeasibility. Original
cost bounds must never be computed from the Phase-I objective.

\section{Branching, incumbents and node bounds}
\label{sec:bp_branching}

At node $N$, let $\mathcal Q_j^N$ contain local patterns consistent with all
inherited decisions. Column generation and bound certification are repeated
over these node-specific sets. Pricing only at the root is insufficient.

\subsection{Branching in the original variable space}

For a fractional master solution, reconstruct all original binaries through
\eqref{eq:bp_projection}. Prefer an assignment component
$0<\bar x_{ij}<1$ and create the children
\begin{equation}
 x_{ij}=0\qquad\text{and}\qquad x_{ij}=1.
 \label{eq:bp_x_branch}
\end{equation}
Each child excludes incompatible existing columns and imposes the same fixing
in every future pricing solve for tail $j$. Other tail pools are unchanged,
although coverage propagation may yield additional valid fixings.

Integral assignments alone do not establish an integral scheduling state.
If the assignment components are integral but a maintenance component is
fractional, branch on an original binary component $v_{jq}$, for example a
$y$, $z$ or hierarchy indicator:
\begin{equation}
 \mathcal Q_j^{N0}=\{v_j\in\mathcal Q_j^N:v_{jq}=0\},\qquad
 \mathcal Q_j^{N1}=\{v_j\in\mathcal Q_j^N:v_{jq}=1\}.
 \label{eq:bp_general_branch}
\end{equation}
These children partition the original integer patterns. Branching only on an
existing $\lambda$ risks allowing future columns to bypass the intended
decision, whereas original-variable fixings are directly enforceable in
pricing.

If every reconstructed original binary is integral, each projected $\bar v_j$
belongs to $\mathcal Q_j$: it is an integer vector satisfying the linear
local constraints by convexity. Together with the master linking rows it is
a compact-model feasible solution, even if some $\lambda$ values remain
fractional. Its objective is preserved. The reconstructed certificate must
nevertheless be checked against all active original constraints before it
updates the incumbent $U$.

\subsection{Pruning and global certification}

A node is closed only by a certified infeasibility result, a valid bound
$L_N\ge U$, or a certified node-optimal integral projection. Discovering an
integral incumbent during incompletely priced optimization does not alone
close the node. A restricted integer master may improve $U$ if its selected
patterns satisfy the original model, but its reported restricted-pool
optimality is not full-problem optimality.

Children inherit the parent's valid lower bound. When new pricing certificates
are available, use the larger of that inherited bound and the new valid
correction. For the open-node set $\mathcal O$ and a finite incumbent, a global
lower bound is
\begin{equation}
 L_{\mathrm{global}}=\min\left\{U,\min_{N\in\mathcal O}L_N\right\}.
 \label{eq:bp_global_lb}
\end{equation}
When no nodes remain, the inner minimum is $+\infty$; a finite incumbent is
then optimal if every closure was justified. Without an incumbent, exhausting
all nodes by certified infeasibility proves infeasibility. At interruption,
open nodes lacking a new pricing certificate remain represented by their
inherited valid bound, potentially $-\infty$.

For finite bounds and nonzero incumbent cost, report
\begin{equation}
 g=100\,\frac{\max\{0,U-L_{\mathrm{global}}\}}{|U|}.
 \label{eq:bp_reported_gap}
\end{equation}
Zero-cost and missing-bound cases require explicit handling rather than a
fabricated finite percentage. All solver calls share a global remaining-time
budget; separately granting the full budget to every pricing call would not
constitute a fixed-budget comparison.

\section{Proposed algorithm}
\label{sec:bp_algorithm}

Algorithm~\ref{alg:bp} describes the intended controller. In exact arithmetic,
finite binary local sets, complete pricing and original-binary branching give
a finite exact search. This is not a practical runtime guarantee. Degenerate
master LPs, expensive pricing and duplicate columns can make convergence slow.
Dual stabilization or heuristic pricing may be useful implementation options,
but certification must use valid duals and complete pricing bounds, not an
unverified stabilized objective.

\begin{algorithm}[htbp]
\caption{Proposed full-ABCD branch-and-price controller}
\label{alg:bp}
\begin{algorithmic}[1]
\State Initialize root, validated seed patterns, $U\gets+\infty$ and open-node queue
\While{the queue is nonempty and global time remains}
 \State Select node $N$; filter pools and pricing models by inherited fixings
 \State Run Phase-I column generation when a feasible master is unavailable
 \If{full Phase-I infeasibility is certified}
  \State Close $N$ and continue
 \EndIf
 \If{Phase I remains unresolved}
  \State Retain $N$ with its inherited bound and mark it unfinished
  \State Continue to another node, or stop if the global budget is exhausted
 \EndIf
 \Repeat
  \State Solve the Phase-II RMP LP and obtain consistent dual prices
  \State Price every tail over $\mathcal Q_j^N$ within the remaining budget
  \State Add new negative-reduced-cost patterns and retain pricing certificates
  \State Update $L_N$ from the pricing correction when its premises hold
 \Until{full LP certification or interruption of column generation}
 \State Validate any reconstructed integer schedule; update $U$ if improved
 \If{$N$ is certified infeasible, bound-prunable, or node-optimal integral}
  \State Close $N$
 \ElsIf{the full node LP is certified and its original projection is fractional}
  \State Branch on an original binary and enqueue both children
 \Else
  \State Retain $N$ as unfinished with a valid bound
 \EndIf
\EndWhile
\State Report incumbent, global bound, status and complete timing/provenance
\end{algorithmic}
\end{algorithm}

\section{Illustration of the restricted-master bound issue}
\label{sec:bp_example}

The following is an algebraic illustration, not an experiment on the ABCD
benchmark. Consider one flight and two tails. Both admit an idle pattern.
The flight costs 100 on tail 1 and 10 on tail 2, and no coupling resource is
binding. Suppose the initial pool contains the flight pattern only for tail 1,
together with both idle patterns. The RMP assigns the flight to tail 1 and
has objective 100, but the actual integer optimum is 10. Thus reporting 100
as a lower bound would be false.

An optimal RMP dual solution is $\pi=100$ and $\alpha_1=\alpha_2=0$.
The missing tail-2 flight pattern has reduced cost $10-100=-90$.
With pricing minima $p_1=0$ and $p_2=-90$, the correction yields
\begin{equation}
 L=100+\min\{0,0\}+\min\{0,-90\}=10.
\end{equation}
Adding that pattern lets the master reach 10. In the integrated model,
station-capacity duals and alternative maintenance placements perform the
same coordinating role, but with much larger local pattern spaces.

\section{Validation and planned computational evaluation}
\label{sec:bp_evaluation}

Before implementation results are admitted to the thesis, the following gates
are required. First, tiny exhaustive local-pattern enumeration must establish
that the priced sets reproduce the selected compact constraints and bounds.
The complete enumerated master LP must match certified column generation;
its integer optimum and branch-and-price must match the compact MILP.
Second, targeted cases must cover shared station bottlenecks, initial A/B
exposure, C/D deadlines, multi-day checks, hierarchy substitution, alternative
maintenance placements, infeasible idle patterns and sparse trigger domains.
Third, branching tests must include integral assignments with fractional
maintenance states, while interruption tests must distinguish feasible pricing
patterns from certified pricing lower bounds.

Matched experiments should then compare the bound-preserving compact LP,
the certified root Dantzig--Wolfe relaxation, the selected strengthened compact
MILP and the full branch-and-price search. Use the same canonical input,
ABCD checks, cap, cost, station accounting, solver version and thread policy.
Capped and unrestricted studies remain separate protocols. Start with small
verified cases before extending the grid to 7-, 14-, 21- and 30-day horizons.

Record the certified lower bound, incumbent, residual gap, proof status,
generated columns, pricing iterations and calls, search-tree nodes, master
and pricing time, construction time and total wall time. Source and input
hashes and all model options are required. Report bound improvement together
with its computational cost; a stronger root bound is not sufficient evidence
of a faster exact algorithm. Memory reductions require actual measurements,
not an inference from the number of master variables alone.

The existing route experiment enumerates routing-only paths and does not
implement pricing or maintenance-aware branch-and-price. Its results therefore
cannot be relabelled as validation of this chapter. Similarly, the 30-day
A-only Gantt illustrates a different solved configuration, not a full-ABCD
branch-and-price schedule.

\section{Review decisions before coding}
\label{sec:bp_review}

The proposed mathematical design should be evaluated on four questions:
whether the local/shared decomposition reproduces every active compact row;
whether C10 trigger accounting is the intended comparison resource; whether
the bound and Phase-I certificates are sufficiently explicit; and whether
exact MILP pricing is an acceptable first implementation despite its potential
cost. Approval of the chapter does not itself establish computational
effectiveness. Implementation and a reproducible experimental study remain
subsequent tasks, and their results must determine the final scope of the
thesis's contribution claims.